\documentclass[10pt,a4paper]{article}
\usepackage[utf8]{inputenc}
\usepackage[T1]{fontenc}
\usepackage{lmodern}
\usepackage[french]{babel}
\usepackage[margin=2.2cm]{geometry}
\usepackage{booktabs}
\usepackage{amsmath}
\usepackage{amssymb}
\usepackage{graphicx}
\usepackage{enumitem}
\setlist{itemsep=1pt,topsep=3pt}
\setlength{\emergencystretch}{3em}
\setlength{\parskip}{3pt}
\setlength{\parindent}{0pt}
\raggedbottom % Évite l’étirement vertical forcé en bas des pages.

% Insère une capture seule depuis le dossier captures/.
\newcommand{\capture}[1]{%
  \IfFileExists{captures/#1}{%
    \par\medskip
    \begin{center}
      \includegraphics[width=0.85\linewidth,keepaspectratio]{captures/#1}
    \end{center}
    \medskip
  }{}%
}

% Insère côte à côte la sortie de la commande et le graphique.
\newcommand{\capturesDeux}[2]{%
  \par\medskip
  \begin{center}
    \begin{minipage}[t]{0.48\linewidth}
      \centering
      \IfFileExists{captures/#1}{%
        \includegraphics[width=\linewidth,keepaspectratio]{captures/#1}\\
        {\footnotesize\textit{Sortie}}
      }{}
    \end{minipage}\hfill
    \begin{minipage}[t]{0.48\linewidth}
      \centering
      \IfFileExists{captures/#2}{%
        \includegraphics[width=\linewidth,keepaspectratio]{captures/#2}\\
        {\footnotesize\textit{Graphe}}
      }{}
    \end{minipage}
  \end{center}
  \medskip
}

\title{Journal de bord - Projet 2, partie 1\\HAI716I : tournoi du dilemme du prisonnier itéré}
\author{Zéïnab Aly CAMARA\\{\small Travail réalisé en binôme avec Neha HAJEE}}
\date{}

\begin{document}
\maketitle

%==============================================================================
\section{Présentation du projet et organisation}
%==============================================================================

Le sujet nous demande de construire un tournoi du dilemme du prisonnier itéré, dans lequel plusieurs stratégies s'affrontent
sur plusieurs parties. Il ne s'agit pas seulement de programmer les stratégies : nous devons aussi observer leurs
comportements, comparer leurs scores et expliquer pourquoi certaines stratégies réussissent mieux que d'autres selon les
règles du tournoi.

Même si le sujet présente le travail comme individuel, nous avons été autorisées à travailler en binôme, à condition
que chacune rédige son propre journal de bord. Nous avons donc réalisé le projet ensemble, tout en conservant chacune notre
journal. Pour nous organiser, nous avons utilisé les \textit{issues} GitHub afin de découper le projet en tâches et de suivre
leur avancement, ce qui nous a permis de travailler en parallèle et de garder une vue d'ensemble.

Toutes les idées n'étaient cependant pas prévues dès le départ dans les issues. Au fil des essais, nous avons ajouté des outils
de comparaison, notamment \texttt{compare\_auto\_jeu.py}, afin de répondre plus précisément à certaines questions soulevées
par les premiers résultats. Le projet a donc évolué avec nos observations et pas uniquement en suivant une liste de tâches
établie au début.

Le moteur du tournoi a été développé par ma binôme. Sa fonction \texttt{jouer\_partie} réinitialise les deux stratégies, leur
demande leur coup à chaque tour en ne leur transmettant que des copies des deux historiques (les coups sont simultanés), puis
ajoute les gains lus dans la matrice de gains et renvoie les deux scores. De mon côté, je me suis notamment occupée des
stratégies, de leurs tests et de mes contributions aux étapes suivantes, en particulier \textbf{GENEROUS-TFT},
\textbf{TIT-FOR-TWO-TATS}, la sélection sur plusieurs générations et les matrices de gains alternatives. Nous avons également
confronté nos choix et nos résultats pour garder une implémentation cohérente.

\subsection*{Difficultés rencontrées pendant l'organisation}
Avant de commencer la programmation, j'ai repris le cours sur les interactions multi-agents. J'ai notamment relu et résumé le
CM4 pour mieux comprendre les stratégies et les notions qui se trouvent derrière le dilemme du prisonnier itéré. Cette reprise
m'a été utile : je voulais comprendre les comportements que nous allions programmer et pas seulement traduire des règles en
Python.

J'ai également rencontré un problème matériel avec mon ordinateur. J'ai dû changer de machine, configurer un nouvel
environnement de travail et transférer le projet ainsi que les éléments nécessaires pour reprendre le développement. Cette
étape m'a fait perdre du temps au milieu du travail, mais j'ai pu remettre en place mon environnement et continuer le projet.

%==============================================================================
\section{Protocole expérimental}
%==============================================================================

\subsection{Paramètres imposés par le sujet}
Pour le tournoi de base, nous avons repris les paramètres indiqués dans l'énoncé :
\begin{center}
\small
\begin{tabular}{@{}lc@{}}
\toprule
Paramètre & Valeur \\
\midrule
Nombre de tours par partie & 200 \\
Nombre de parties par paire de stratégies & 5 \\
Graine aléatoire & 42 \\
Gains (C,C) / (C,D) / (D,C) / (D,D) & 3/3 \quad 0/5 \quad 5/0 \quad 2/2 \\
\bottomrule
\end{tabular}
\end{center}

Les coups sont simultanés : au moment de choisir C ou D, une stratégie ne connaît pas encore le coup que l'autre joueur va jouer
au même tour. Elle peut en revanche utiliser l'historique des tours précédents.

Nous avons fixé la graine aléatoire à 42 afin de pouvoir reproduire les expériences. Les stratégies qui utilisent du hasard
reçoivent le générateur aléatoire du moteur (\texttt{rng}) au lieu d'utiliser directement le générateur global de Python. Cela
permet de contrôler les tirages avec la graine choisie.

\subsection{Le choix concernant l'auto-jeu}
Nous avons choisi, pour le tournoi de base, qu'une stratégie ne joue pas contre elle-même. Elle peut néanmoins rencontrer une
autre instance de la même stratégie : il s'agit alors de deux joueurs distincts qui appliquent la même Règle.

Ce choix vise à ne pas avantager automatiquement les stratégies coopératives. Par exemple, deux copies de TIT-FOR-TAT coopèrent
à chaque tour et obtiennent chacune 3 points par tour, soit 600 points sur une partie de 200 tours. Deux copies de ALL-D
obtiennent chacune 2 points par tour, soit 400 points. Ce choix correspond également au mini-tournoi présenté dans le cours, où
un joueur ne s'affronte pas lui-même. Nous avons toutefois voulu mesurer l'effet de cette décision au lieu de supposer qu'elle
était sans conséquence. C'est pour cette raison que nous avons ajouté une comparaison entre un tournoi sans auto-jeu et un
tournoi avec auto-jeu.

\subsection{Quelques règles d'implémentation}
Nous avons représenté les coups par les caractères \texttt{"C"} et \texttt{"D"}. Les historiques sont transmis aux stratégies à
chaque tour : une stratégie peut donc prendre sa décision à partir de ses propres coups et de ceux de son adversaire. Si deux
stratégies ont exactement le même score, la première dans l'ordre alphabétique est classée devant. Nous avons également fixé
deux interprétations pour éviter les ambiguïtés :
\begin{itemize}
  \item \textbf{TESTER} observe la réponse de l'adversaire à sa trahison initiale. Comme les coups sont simultanés, elle ne
  peut pas connaître la réponse au moment de jouer le deuxième tour ; elle prend donc sa décision à partir du coup adverse
  observé au deuxième tour.
  \item \textbf{JOSS} suit le comportement de TIT-FOR-TAT, mais remplace un coup C par D avec une probabilité de 10\,\%. Nous
  avons retenu un tirage aléatoire plutôt qu'une trahison mécanique tous les dix tours, car un rythme fixe serait plus facile à
  anticiper.
\end{itemize}

%==============================================================================
\section{Les stratégies}
%==============================================================================
Cette partie présente les règles utilisées, les comportements que nous pouvions attendre avant de comparer les résultats (des
hypothèses, qui ne garantissent pas qu'une stratégie obtiendra le meilleur score dans tous les tournois) et ce que les scores
par paire ont montré ensuite. Les scores cités sont les scores moyens par partie de 200 tours (graine 42, 5 parties), regroupés
dans le tableau~\ref{tab:matrice}. Les cinq premières stratégies viennent du cours (diapo~64).

\subsection{ALL-D}
\textbf{Règle :} ALL-D joue toujours D, quel que soit l'historique de la partie.

\textbf{Comportement attendu :} Cette stratégie peut obtenir un bon score contre les stratégies qui commencent par coopérer ou
qui pardonnent facilement, puisqu'elle profite de leurs coups C. En revanche, contre une autre ALL-D, les deux joueurs
obtiennent seulement 2 points par tour. Elle ne cherche jamais à établir une coopération.

\textbf{Observé.} Elle ne perd aucun duel, mais souvent de peu (403 contre 398 face à TIT-FOR-TAT). Elle écrase RANDOM (688
contre 208) et ne fait que 400 contre une copie.

\subsection{RANDOM}
\textbf{Règle :} RANDOM choisit C ou D de manière aléatoire, avec une probabilité de 50\,\% pour chaque coup.

\textbf{Comportement attendu :} Son comportement n'est pas adapté à un adversaire particulier : elle ne cherche ni à coopérer
durablement ni à punir une trahison. Elle peut obtenir des résultats variables selon les tirages et selon la manière dont les
autres stratégies réagissent à ses coups.

\textbf{Observé.} Elle perd lourdement contre ALL-D (208 contre 688) et contre JOSS (480 contre 521) et fait jeu égal avec
TIT-FOR-TAT (500 contre 497).

\subsection{TIT-FOR-TAT}
\textbf{Règle :} TIT-FOR-TAT commence par C, puis reproduit au tour suivant le dernier coup joué par son adversaire. C'est la
stratégie gagnante du tournoi d'Axelrod de 1980.

\textbf{Comportement attendu :} Elle commence par faire confiance, répond à une trahison par une trahison, mais revient à la
coopération dès que l'adversaire coopère de nouveau. Elle devrait donc bien fonctionner avec des stratégies coopératives ou
capables de revenir à la coopération. Elle risque en revanche d'être exploitée par une stratégie qui trahit dès le début.

\textbf{Observé.} Elle ne gagne aucun duel nettement (398 contre 403 face à ALL-D, 414 contre 419 face à JOSS), mais elle fait
600 contre une copie. Contre TESTER, elle fait 500 contre 500 : une alternance où chacun reçoit tour à tour 5 et 0.

\subsection{TESTER}
\textbf{Règle :} TESTER commence par D afin de tester la réaction de son adversaire. Si celui-ci répond D au deuxième tour,
TESTER adopte ensuite le comportement de TIT-FOR-TAT. Dans le cas contraire, elle répète le motif C, C, D.

\textbf{Comportement attendu :} Cette stratégie tente de repérer les adversaires qui se laissent exploiter. Elle peut profiter
des stratégies qui ne ripostent pas, mais son test initial peut être coûteux face à une stratégie qui répond immédiatement à la
trahison.

\textbf{Observé.} Elle perd de peu contre ALL-D (398 contre 403) et fait 500 contre 500 face à TIT-FOR-TAT. Dans la population
élargie, elle exploite fortement TIT-FOR-TWO-TATS (734 contre 399).

\subsection{JOSS}
\textbf{Règle :} JOSS suit TIT-FOR-TAT, mais lorsqu'elle devrait jouer C, elle joue D avec une probabilité de 10\,\%. Lorsqu'elle
devrait déjà jouer D, elle conserve D.

\textbf{Comportement attendu :} Ses trahisons occasionnelles peuvent lui permettre de profiter d'adversaires indulgents. En
revanche, contre une stratégie qui réagit immédiatement comme TIT-FOR-TAT, une trahison aléatoire peut déclencher une période
de représailles.

\textbf{Observé.} Elle bat légèrement TIT-FOR-TAT (419 contre 414) et RANDOM (521 contre 480), mais ne fait que 408 contre une
copie, loin des 600 de TIT-FOR-TAT : ses trahisons déclenchent des échos de représailles.

\subsection{GENEROUS-TFT - stratégie ajoutée par mes soins}
\textbf{Règle :} GENEROUS-TFT commence par C. Si l'adversaire a joué C au tour précédent, elle joue C. Si l'adversaire a joué D,
elle lui pardonne avec une probabilité de 1/3 ; sinon, elle riposte par D. Dans le code, la classe Python s'appelle
\texttt{GenerousTitForTat} et son nom dans le classement est \texttt{GENEROUS-TFT}.

\textbf{Origine et choix du paramètre.} Cette stratégie reprend l'idée de TIT-FOR-TAT généreux évoquée dans le cours
(Nowak et Sigmund), notamment comme moyen de limiter les cycles de représailles. La probabilité de pardon est calculée à partir
des gains du dilemme avec la formule suivante, où $T=5$, $R=3$, $P=2$ et $S=0$ :
\[
g = \min\left(1-\frac{T-R}{R-S},\frac{R-P}{T-P}\right)
  = \min\left(\frac{1}{3},\frac{1}{3}\right)=\frac{1}{3}.
\]

\textbf{Comportement attendu :} Par rapport à TIT-FOR-TAT, cette stratégie est plus indulgente : une trahison adverse ne provoque
pas nécessairement une riposte. Nous voulions ainsi étudier si le pardon pouvait aider à éviter les enchaînements de
représailles, notamment face à des stratégies qui introduisent parfois une trahison.

\textbf{Observé (population élargie).} Elle fait 600 contre TIT-FOR-TAT, TIT-FOR-TWO-TATS et une copie, mais elle est exploitée
par ALL-D (266 contre 600), RANDOM (422 contre 605) et TESTER (553 contre 600).

\subsection{TIT-FOR-TWO-TATS - stratégie ajoutée par mes soins}
\textbf{Règle :} TIT-FOR-TWO-TATS commence par C et ne riposte que si l'adversaire a joué D deux fois de suite. Une trahison
isolée ne suffit donc pas à déclencher une punition. C'est une variante classique, plus tolérante, de TIT-FOR-TAT.

\textbf{Comportement attendu :} Elle peut éviter de réagir trop fortement à une trahison ponctuelle, mais cette tolérance peut
aussi être exploitée par des adversaires qui trahissent de temps en temps sans jouer deux D consécutifs.

\textbf{Observé (population élargie).} Elle fait 600 contre TIT-FOR-TAT, GENEROUS-TFT et une copie, mais elle est très exploitée
par TESTER (399 contre 734), RANDOM (403 contre 657) et JOSS (555 contre 630).

%==============================================================================
\section{Étape 1 - Tournoi de base}
%==============================================================================

\subsection{Résultats sans auto-jeu}
Nous avons lancé le tournoi avec les cinq stratégies du cours, une graine fixée à 42, cinq parties de 200 tours par
confrontation et sans auto-jeu. Le classement obtenu est le suivant :
\begin{center}
\small
\begin{tabular}{@{}rlr@{}}
\toprule
Rang & Stratégie & Score total \\
\midrule
1 & ALL-D & \textbf{9\,482} \\
2 & TIT-FOR-TAT & 9\,056 \\
3 & TESTER & 8\,976 \\
4 & JOSS & 8\,728 \\
5 & RANDOM & 8\,477 \\
\bottomrule
\end{tabular}
\end{center}

\textbf{La stratégie gagnante est ALL-D.} Ce résultat nous a surpris, car nous pensions que TIT-FOR-TAT gagnerait grâce à sa
capacité à s'adapter au comportement de l'adversaire. Les scores par paire (tableau~\ref{tab:matrice}) montrent d'où vient
l'avantage de ALL-D.

\begin{table}[htbp]
\centering
\small
\begin{tabular}{@{}lrrrrr@{}}
\toprule
  & ALL-D & RANDOM & TFT & TESTER & JOSS \\
\midrule
ALL-D & \textit{400} & 688 & 403 & 403 & 402 \\
RANDOM & 208 & \textit{499} & 500 & 532 & 480 \\
TIT-FOR-TAT & 398 & 497 & \textit{600} & 500 & 414 \\
TESTER & 398 & 483 & 500 & \textit{533} & 405 \\
JOSS & 398 & 521 & 419 & 405 & \textit{408} \\
\bottomrule
\end{tabular}
\caption{Score moyen par partie de la stratégie en ligne contre celle en colonne (graine 42). La diagonale, en italique, est
le score contre une copie : elle n'entre pas dans le tournoi sans auto-jeu.}
\label{tab:matrice}
\end{table}

Contre les quatre autres, ALL-D totalise $688+403+403+402=1\,896$ points et TIT-FOR-TAT $398+497+500+414=1\,809$. L'écart de 87
se décompose ainsi : $+191$ pour ALL-D contre RANDOM (688 contre 497), $-97$ contre TESTER (403 contre 500), $-12$ contre JOSS
(402 contre 414) et $+5$ dans leur duel. ALL-D gagne donc surtout parce qu'elle exploite RANDOM, la seule stratégie de la
population qui se laisse exploiter sans réagir. Face aux autres, elle ne gagne presque rien et TIT-FOR-TAT la devance même
contre TESTER. Ce résultat montre que la stratégie qui semble la plus raisonnable dans une interaction particulière n'est pas
forcément celle qui obtient le meilleur score total dans un tournoi donné.

\subsection{Comparaison avec auto-jeu}
Pour mesurer l'effet de notre choix de ne pas inclure l'auto-jeu, nous avons utilisé \texttt{compare\_auto\_jeu.py}. Le
programme compare les classements obtenus avec et sans confrontation d'une stratégie contre une copie de sa propre Règle :
\begin{center}
\small
\begin{tabular}{@{}lrrrr@{}}
\toprule
Stratégie & Sans : rang & Sans : score & Avec : rang & Avec : score \\
\midrule
ALL-D & 1 & 9\,482 & 3 & 11\,482 \\
TIT-FOR-TAT & 2 & 9\,056 & 1 & \textbf{12\,042} \\
TESTER & 3 & 8\,976 & 2 & 11\,594 \\
JOSS & 4 & 8\,728 & 5 & 10\,753 \\
RANDOM & 5 & 8\,477 & 4 & 11\,118 \\
\bottomrule
\end{tabular}
\end{center}

Avec l'auto-jeu, \textbf{TIT-FOR-TAT devient la gagnante}. Ses parties contre une copie lui rapportent environ
$5\times600=3\,000$ points, contre $5\times400=2\,000$ pour ALL-D, ce qui suffit à renverser l'ordre. Ce n'est pas un effet de
la graine : sur 30 graines, ALL-D gagne 30 fois sans auto-jeu et TIT-FOR-TAT gagne 30 fois avec. Cette comparaison confirme que le
classement dépend aussi du protocole expérimental. Il ne suffit donc pas de dire qu'une stratégie est \og la meilleure \fg{}
sans préciser la population et les règles du tournoi.

\subsection{Réponse à la question du sujet : TIT-FOR-TAT gagne-t-elle chez nous ?}
Dans notre tournoi de base \textbf{sans auto-jeu, TIT-FOR-TAT ne gagne pas} : elle termine deuxième, derrière ALL-D. En revanche,
elle termine première lorsque l'auto-jeu est autorisé. Nous ne pouvons donc pas attribuer la victoire de TIT-FOR-TAT à son seul
comportement : le résultat dépend aussi des stratégies présentes (la population est petite et peu coopérative, seule
TIT-FOR-TAT coopère vraiment) et de la possibilité de rencontrer une copie de soi-même.

%==============================================================================
\section{Étape 2 - Ajout de nouvelles stratégies}
%==============================================================================
Pour élargir la population, nous avons ajouté deux stratégies qui ne figuraient pas parmi les cinq stratégies du cours :
GENEROUS-TFT, qui peut pardonner une trahison avec une probabilité de 1/3 et TIT-FOR-TWO-TATS, qui attend deux trahisons
consécutives avant de riposter. Ces deux règles permettent de comparer deux formes d'indulgence.

\subsection{Classement de la population élargie sans auto-jeu}
Avec les sept stratégies présentes dans la population étendue du code, nous obtenons :
\begin{center}
\small
\begin{tabular}{@{}rlr@{}}
\toprule
Rang & Stratégie & Score total \\
\midrule
1 & TESTER & \textbf{15\,490} \\
2 & TIT-FOR-TAT & 15\,087 \\
3 & RANDOM & 15\,068 \\
4 & JOSS & 15\,026 \\
5 & GENEROUS-TFT & 14\,706 \\
6 & TIT-FOR-TWO-TATS & 14\,673 \\
7 & ALL-D & 14\,514 \\
\bottomrule
\end{tabular}
\end{center}

\textbf{La stratégie gagnante est TESTER.} L'ajout de stratégies modifie le classement : ALL-D, première dans le tournoi de
base, se retrouve dernière. Les scores par paire l'expliquent : ALL-D ne dépasse 400 points que contre RANDOM (688) et
GENEROUS-TFT (600), alors que les autres font 600 contre les stratégies indulgentes et que TESTER exploite TIT-FOR-TWO-TATS
(734 contre 399). Ajouter des stratégies coopératives profite donc davantage à celles qui coopèrent qu'à ALL-D.

Les deux stratégies que j'ai ajoutées ne prennent pas les premières places. Cela ne signifie pas que leur principe est inutile :
elles ont été conçues pour tester des comportements particuliers, notamment le pardon et la tolérance et pas pour garantir une
première place dans toutes les populations. Le résultat est robuste (sur 30 graines, TESTER gagne 30 fois), mais les écarts en
tête restent faibles : ce classement vaut pour les paramètres et les stratégies de cette expérience.

\subsection{Comparaison avec et sans auto-jeu dans la population élargie}
\begin{center}
\small
\begin{tabular}{@{}lrrrr@{}}
\toprule
Stratégie & Sans : rang & Sans : score & Avec : rang & Avec : score \\
\midrule
TESTER & 1 & 15\,490 & 1 & \textbf{18\,175} \\
TIT-FOR-TAT & 2 & 15\,087 & 2 & 18\,061 \\
RANDOM & 3 & 15\,068 & 5 & 17\,527 \\
JOSS & 4 & 15\,026 & 6 & 16\,970 \\
GENEROUS-TFT & 5 & 14\,706 & 3 & 17\,847 \\
TIT-FOR-TWO-TATS & 6 & 14\,673 & 4 & 17\,761 \\
ALL-D & 7 & 14\,514 & 7 & 16\,514 \\
\bottomrule
\end{tabular}
\end{center}
Dans cette population élargie, TESTER reste première avec ou sans auto-jeu. L'auto-jeu modifie néanmoins les rangs de plusieurs
stratégies : GENEROUS-TFT passe notamment de la cinquième à la troisième place, car elle fait 600 contre une copie.

%==============================================================================
\section{Étape 3 - Modifications du tournoi}
%==============================================================================
Pour cette étape, nous avons étudié deux modifications prévues par le sujet : la sélection sur plusieurs générations et le
remplacement de la matrice de gains du dilemme par d'autres matrices. Nous avons choisi la sélection parce qu'un tournoi
additionne des scores une seule fois, alors que le succès d'une stratégie dépend de qui l'entoure : la sélection demande si
elle survit quand sa part dépend de ses résultats. Nous avons choisi d'autres matrices (chasse au cerf, poule mouillée, vues dans
le cours) pour vérifier si nos conclusions dépendent du dilemme lui-même.

\subsection{Sélection sur plusieurs générations}
\textbf{Principe.} Nous avons implémenté une dynamique inspirée de la dynamique de réplication. Au départ, toutes les
stratégies ont la même proportion. La fitness d'une stratégie dépend de son score moyen face à la population présente et
sa proportion augmente si elle dépasse la fitness moyenne :
\[
p_i' = p_i \times \frac{f_i}{\bar f},
\]
où \(p_i\) est la proportion de la stratégie \(i\), \(f_i\) sa fitness et \(\bar f\) la fitness moyenne. La matrice des scores
moyens est calculée une fois avec la graine, puis réutilisée. Les proportions restent normalisées et il n'y a pas de mutation :
une proportion exactement nulle ne peut pas réapparaître.

\textbf{Population du cours} (100 générations, matrice du dilemme, graine 42) : TIT-FOR-TAT passe de 0,200 à \textbf{1,000} et
les quatre autres (TESTER, RANDOM, JOSS, ALL-D) tombent à 0,000.

Ce résultat est différent du tournoi sans auto-jeu, où ALL-D arrivait première. L'explication est dans la diagonale de la
matrice : à la génération 0, la fitness de chaque stratégie est la moyenne de sa ligne du tableau~\ref{tab:matrice}, copie
comprise (TIT-FOR-TAT 482, TESTER 464, ALL-D 459, RANDOM 444, JOSS 430). C'est exactement l'ordre du tournoi \emph{avec}
auto-jeu : la sélection fait rencontrer des copies, comme ce tournoi et non comme le tournoi de référence. Ensuite, les
stratégies faibles disparaissent et retirent leurs proies à ALL-D, qui ne fait que 400 contre une copie. Ce résultat dépend
aussi des proportions initiales et de l'absence de mutation.

\textbf{Population élargie} (sept stratégies, 0,143 chacune au départ) :
\begin{center}
\small
\begin{tabular}{@{}lr@{}}
\toprule
Stratégie & Proportion à la génération 100 \\
\midrule
TIT-FOR-TAT & \textbf{0,591} \\
GENEROUS-TFT & \textbf{0,382} \\
TIT-FOR-TWO-TATS & 0,022 \\
RANDOM & 0,003 \\
TESTER & 0,002 \\
JOSS, ALL-D & 0,000 \\
\bottomrule
\end{tabular}
\end{center}
TIT-FOR-TAT reste la plus représentée, mais ne prend plus toute la population. GENEROUS-TFT se maintient à environ 38\,\%, bien
mieux que TIT-FOR-TWO-TATS : les deux font 600 l'une contre l'autre et contre elles-mêmes, mais TIT-FOR-TWO-TATS est exploitée
par TESTER, RANDOM et JOSS. C'est une observation liée à notre matrice de scores et à notre modèle de sélection et non une
propriété universelle de ces stratégies.

\subsection{Modification des matrices de gains}
Les matrices disponibles dans le code sont : le dilemme du cours (3/3, 0/5, 5/0, 2/2), le dilemme avec une punition plus
sévère (D,D = 1/1), la chasse au cerf (4/4, 0/3, 3/0, 1/1) et la poule mouillée (3/3, 1/4, 4/1, 0/0). Les deux jeux
alternatifs étudiés ici sont :
\begin{itemize}
  \item \textbf{La chasse au cerf} : la coopération mutuelle est particulièrement intéressante, mais elle suppose que les
  joueurs se fassent confiance. La matrice comporte deux équilibres de Nash, correspondant à la coopération mutuelle et à la
  défection mutuelle.
  \item \textbf{La poule mouillée} : le pire résultat pour les deux joueurs est la défection mutuelle. Il peut donc être
  risqué de toujours choisir D, contrairement à une interprétation trop simplifiée du dilemme du prisonnier.
\end{itemize}
Les valeurs numériques de la poule mouillée ont été choisies dans le code pour respecter l'ordre de préférence indiqué dans le
cours ; elles ne sont pas présentées comme des valeurs imposées par le support.

\textbf{Chasse au cerf.} Avec la population du cours, la graine 42 et 100 générations, TIT-FOR-TAT atteint une proportion
finale de 1,000, les quatre autres 0,000. Ici, trahir un coopérateur rapporte 3, moins que la coopération mutuelle (4) :
l'exploitation ne paie plus. Deux copies de TIT-FOR-TAT font 800 points, deux copies de ALL-D seulement 200. TIT-FOR-TAT est
donc la gagnante de cette expérience de sélection.

\textbf{Poule mouillée.} TIT-FOR-TAT atteint également 1,000. Deux copies de ALL-D font 0 point, car la défection mutuelle vaut
0/0 et ALL-D ne gagne que 4 points contre TIT-FOR-TAT (un seul tour d'avance avant que la défection mutuelle s'installe).
JOSS souffre aussi de ses échos de représailles (32 points contre une copie). Une stratégie qui trahit systématiquement ne
bénéficie donc pas de son comportement, alors que la coopération conditionnelle de TIT-FOR-TAT reste avantageuse.

Ces résultats ne signifient pas que TIT-FOR-TAT gagnerait dans toutes les variantes de ces jeux : ils décrivent les résultats de
notre implémentation, avec les stratégies, les paramètres et les matrices de gains retenus.

%==============================================================================
\section{Vérifications et tests}
%==============================================================================
Nous avons écrit des tests pour vérifier plusieurs comportements importants :
\begin{itemize}
  \item les règles des stratégies, notamment les premiers coups et les réactions aux historiques et les valeurs du cours
  (TIT-FOR-TAT contre ALL-D sur 10 tours donne 18 contre 23, le mini-tournoi donne 48 / 48 / 46) ;
  \item le comportement probabiliste de RANDOM, JOSS et GENEROUS-TFT sur un grand nombre de tours ;
  \item la reproductibilité des stratégies aléatoires avec une même graine ;
  \item le fonctionnement de la comparaison avec et sans auto-jeu ;
  \item les propriétés de la sélection, notamment la somme des proportions et l'absence de réapparition d'une proportion nulle.
\end{itemize}
La commande \texttt{python3 -m pytest -q} affiche \textbf{29 tests réussis} (capture en section~\ref{sec:captures}).

%==============================================================================
\section{Difficultés, choix et limites}
%==============================================================================

\subsection{Comprendre les stratégies avant de les programmer}
La reprise du CM4 m'a aidée à distinguer les stratégies qui coopèrent, celles qui trahissent systématiquement et celles qui
adaptent leur comportement à l'historique. TESTER et JOSS m'ont notamment demandé de bien réfléchir à la notion de réaction à
l'adversaire, car leurs décisions dépendent de ce qui s'est passé pendant la partie.

\subsection{Interpréter les résultats plutôt que regarder uniquement le classement}
L'un des enseignements importants du projet est qu'un classement ne suffit pas à expliquer le comportement des stratégies.
ALL-D gagne le tournoi de base sans auto-jeu, TIT-FOR-TAT gagne lorsque l'auto-jeu est autorisé, TESTER gagne dans la
population élargie sans auto-jeu, tandis que TIT-FOR-TAT domine la sélection sur plusieurs générations dans les expériences
réalisées. Ces résultats correspondent à des protocoles différents : il faut toujours préciser les paramètres et les règles
avant d'interpréter un gagnant.

\subsection{Limites de nos expériences}
Nos résultats sont obtenus avec une graine déterminée (dont nous avons vérifié la stabilité sur 30 graines pour les gagnants),
un nombre précis de parties et une population limitée. La sélection repose sur une matrice de scores calculée au départ et ne
prévoit pas de mutation. Les résultats décrivent donc ce modèle particulier et non toutes les formes possibles d'évolution
d'une population.

%==============================================================================
\section{Bilan personnel}
%==============================================================================
Ce projet nous a permis de passer des notions vues en cours à une expérimentation concrète. Nous avons pu observer que le
comportement d'une stratégie ne peut pas être évalué indépendamment de ses adversaires et des règles du tournoi.

Le résultat qui m'a le plus marquée est le changement de gagnant selon le protocole : ALL-D gagne le tournoi de base sans
auto-jeu, TIT-FOR-TAT gagne lorsque l'auto-jeu est ajouté et TESTER arrive en tête lorsque la population est élargie. La
sélection sur plusieurs générations donne encore une autre lecture, en faisant progressivement dominer TIT-FOR-TAT dans les
expériences réalisées.

L'ajout de GENEROUS-TFT et de TIT-FOR-TWO-TATS m'a également permis d'explorer deux façons de gérer les trahisons : pardonner
avec une certaine probabilité ou attendre plusieurs trahisons avant de riposter. Les résultats de la sélection élargie montrent
notamment que GENEROUS-TFT conserve une place importante dans la population.

Enfin, ce travail m'a appris à ne pas m'arrêter au premier résultat obtenu. Il faut comprendre les règles de l'expérience,
vérifier que le programme les applique correctement, comparer plusieurs configurations et expliquer les écarts observés. C'est
cette démarche expérimentale que je retiens principalement de ce projet.

%==============================================================================
\clearpage % La partie des captures commence sur une nouvelle page.
\section{Captures d'écran}\label{sec:captures}
%==============================================================================
Chaque capture correspond à une exécution réelle. La commande de lancement est
présentée, puis la ou les captures correspondantes sont insérées juste en dessous.
Les images sont placées dans le dossier \texttt{captures/}, à côté de ce fichier
LaTeX.

\medskip
\textbf{1. Tournoi de base et comparaison avec/sans auto-jeu}
{\footnotesize\begin{verbatim}
python3 compare_auto_jeu.py --seed 42
\end{verbatim}}
\capture{1_COMPARE.png}

\medskip
\textbf{2. Tournoi avec la population élargie}
{\footnotesize\begin{verbatim}
python3 compare_auto_jeu.py --seed 42 --population etendue
\end{verbatim}}
\capture{2_POPULTATION_ELARGIE.png}

\medskip
\textbf{3. Sélection avec la population du cours}
{\footnotesize\begin{verbatim}
python3 selection.py --seed 42 --generations 100 --csv sel_cours.csv --png sel_cours.png
\end{verbatim}}
\capturesDeux{3_SELECTION_POP_COURS.png}{3_GRAPHIQUE_POP_COURS.png}

\medskip
\textbf{4. Sélection avec la population élargie}
{\footnotesize\begin{verbatim}
python3 selection.py --seed 42 --generations 100 --population etendue --png sel_etendue.png
\end{verbatim}}
\capturesDeux{4_SELECTION_POP_ELARGIE.png}{4_GRAP_SELECTION_POP_ELARGIE.png}

\medskip
\textbf{5. Sélection avec la chasse au cerf}
{\footnotesize\begin{verbatim}
python3 selection.py --seed 42 --generations 100 --gains cerf --png sel_cerf.png
\end{verbatim}}
\capturesDeux{5_CHASSE_CERF.png}{5_GRAPH_CHASSE_CERF.png}

\medskip
\textbf{6. Sélection avec la poule mouillée}
{\footnotesize\begin{verbatim}
python3 selection.py --seed 42 --generations 100 --gains poule --png sel_poule.png
\end{verbatim}}
\capturesDeux{6_POULE_MOUILLEE.png}{6_GRAPH_POULE_MOUILLEE.png}

\medskip
\textbf{7. Tests automatisés}
{\footnotesize\begin{verbatim}
python3 -m pytest -q
\end{verbatim}}
\capture{7_TESTS_AUTOMATISES.png}

\end{document}
